\section{Blacbox PIT for Depth-4 Circuits} \label{sec:pf-thm2}

We will give the proof of \autoref{thm:thm2} in this section. Before the details, we will state a few important definitions and lemmas from \cite{agrawal2016jacobian} to be referenced later.

\begin{definition}[Transcendence Degree]
    Polynomials $T_1,\hdots,T_m$ are called {\em algebraically dependent} if there exists a nonzero {\em annihilator} $A$ s.t.~$A(T_1,\hdots,T_m) = 0$. {\em Transcendence degree} is the size of the largest subset $S\subseteq \{T_1,\hdots,T_m\}$ that is algebraically independent. Then $S$ is called a {\em transcendence basis}.
\end{definition}

\begin{definition}[Faithful homomorphism]
    A homomorphism $\Phi : \F[\x] \to \F[\y]$ is faithful for $\T$ if $\tdeg_{\F}(\T) = \tdeg_{\F}(\Phi(\T))$.
\end{definition}
    
The reason for interest in faithful maps is due its usefulness in preserving the identity as shown in the following fact.

\begin{fact}[Theorem 2.4 \cite{agrawal2016jacobian}]\label{fact:faithful-is-useful}
    For any $C \in \F[y_1,\hdots,y_m]$, $C(\T) = 0 \iff C(\Phi(\T))=0$. 
\end{fact}

Here is an important criterion about the jacobian matrix which basically shows that it \emph{preserves} algabraic independence.
    
\begin{fact}[Jacobian criterion] \label{fact:jacobian-criterion}
    Let ${\bf f} \subset \F[\x]$ be a finite set of polynomials of degree at most $d$, and $\tdeg_{\F} ({\bf f}) \le r$. If char$(\F) = 0$, or char$(\F) > d^r$, then $\tdeg_{\F} ({\bf f}) =
    \rk_{\F(x)} \calJ_{\x}({\bf f})$.
\end{fact}

Jacobian criterion together with faithful maps give a recipe to design a map which drastically reduces number of variables, if trdeg is small.

\begin{lemma}[Lemma 2.7 \cite{agrawal2016jacobian}] \label{lem:faithful}
    Let $\T \in \F[\x]$ be a finite set of polynomials
    of degree at most $d$ and $\tdeg_{\F}(\T) \le r$, and char(F)=$0$ or $> d^r$. Let $\Psi' : \F[\x] \longrightarrow \F[z]$ such that $\rk_{\F(\x)} \calJ_{\x}(\T) = \rk_{\F(z)} \Psi'(\calJ_{\x}(\T))$. 
    
    Then, the map $\Phi : \F[\x] \longrightarrow \F[z,t,\y]$, such that $x_i \mapsto (\sum_{j\in [r]} y_j t^{ij}) + \Psi'(x_i)$, is a faithful homomorphism for $\T$.
\end{lemma}

In the next section we will use these tools to prove \autoref{thm:thm2}(b). The proof and calculations for \autoref{thm:thm2}(a) are very similar.

\subsection{PIT for \texorpdfstring{$\SPSP{[k]}{}{[\delta]}$}{SkPSPdelta}} \label{sec:SkPSPdelta}

We solve the PIT for a more general model than $\SPSP{[k]}{}{}$ by solving the following problem.

\begin{problem}\label{prob:prob1-thm2}
 Let $\{T_i \,|\, i \in [m]\}$ be $\PSP{}{[\delta]}$ circuits of (syntactic) degree at most $d$ and size $s$. Let the transcendence degree of $T_i$'s, $\tdeg_{\F}(T_1,\hdots,T_m)= k\ll s$. Further, $C(x_1,\hdots,x_m)$ be a circuit of $(\size + \deg) < s'$. Design a blackbox-PIT algorithm for $C(T_1,\hdots,T_m)$.
\end{problem}

Trivially, $\SPSP{[k]}{}{[\delta]}$ is a very special case of the above setting. Let $\T:=\{T_1,\hdots,T_m\}$. Let $\T_{k}:=\{T_1,\hdots,T_k\}$ be a transcendence basis. For $T_i=\prod_{j} g_{ij}$, we denote the set $L(T_i):=\{g_{ij}\,\mid\,j\}$. 

We want to find an explicit homomorphism $\Psi : \F[\x] \to \F[\x,z]$ s.t.~$\Psi(\calJ_{\x}(\T))$ is of a `nice' form. In the image we fix $\x$ suitably, to get a composed map $\Psi' : \F[\x] \longrightarrow \F[z]$ s.t.~$\rk_{\F(\x)}\calJ_{\x}(\T) = \rk_{\F(z)}\Psi'(\calJ_{\x}(\T))$. Then, we can extend this map to $\Phi : \F[\x] \longrightarrow \F[z,\y,t]$ s.t.~$x_i \mapsto (\sum_{j=1}^k\,y_jt^{ij}) + \Psi'(x_i)$, which is \emph{faithful} \autoref{lem:faithful}. We show that the map $\Phi$ can be efficiently constructed using a scaling and shifting map ($\Psi$) which is eventually fixed by the hitting set ($H'$ defining $\Psi'$) of a $\SWSP{}{}{[\delta]}$ circuit. Overall, $\Phi(f)$ is a $k+2$-variate polynomial for which a trivial hitting set exists.

Wlog, $\calJ_{\x}(\T)$ is full rank with respect to the variable set $\x_{k}=(x_1,\hdots,x_k)$. Thus, by assumption, $J_{{\x}_k}(\T_k) \ne 0$ (for notation, see section \ref{sec:prel}). We want to construct a $\Psi$ s.t.~$\Psi(J_{{\x}_k}(\T_k))$ has an `easier' PIT. We have the following identity \cite[Eqn.~3.1]{agrawal2016jacobian}, from the linearity of the determinant, and the simple observation that $\partial_{x}(T_i) \;=\; T_i \cdot \left(\sum_{j}\,\partial_{x}(g_{ij})/g_{ij}\right)$, where $T_i\,=\,\prod_{j}\,g_{ij} $:
\begin{align}
    J_{{\x}_k}(\T_k)\;=\;\sum_{g_1 \in L(T_1) , \hdots , g_k \in L(T_k)}\;\; \left( \frac{T_1\hdots T_k}{g_1 \hdots g_k}\right) \cdot J_{{\x}_k}(g_1,\hdots,g_k)\;. \label{jacobian-main-eq}
\end{align}

\textbf{The homomorphism $\Psi$.} To ensure the invertibility of all $g \in \bigcup_{i}\,L(T_i)$ we proceed as in section \ref{sec:pf:thm1}. Consider 
\begin{align*}
    h := \prod_{i \in [k]} \prod_{g \in L(T_i)} g = \prod_{i \in [\ell]} g,
\end{align*}
where $g \in \bigcup_{i}\,L(T_i)$ and $\ell \leq k \cdot s$. Note that $\deg h \leq d \cdot k \cdot s$ and $h$ is computable by $\PSP{}{}$ circuit of size $O(s)$. \autoref{lem:prod-sparse-hs} gives the relevant hitting set $\calH \subseteq \F^n$ which contains an evaluation point $\alphavec = (a_1, \dots, a_n)$ such that $h(\alphavec) \neq 0$ implying $g(\alphavec) \neq 0$, for all $g \in \bigcup_{i}\,L(T_i)$. We emphasise that, unlike the previous case, here in the blackbox setting, we {\em do not} have individual access of $g$ to verify for the correct $\alphavec$. Thus, we try out all $\alphavec \in \calH$ to see whichever works. If the input polynomial $f$ is non-zero, then one such $\alphavec$ must exist. This search adds a multiplicative blowup of $\poly(s)$, since the size of $\calH$ is $\poly(s)$. 

Fix an $\alphavec = (a_1, \cdots, a_n) \in \calH$ and define $\Psi: \F[\x] \to \F[\x,z]$ as $x_i \mapsto z \cdot x_i + a_i$.
%
Denote the ring $\ringR[\x]$ where $\ringR:=\F[z]/\langle z^D \rangle$, and $D:=k\cdot(d-1)+1$. Being 1-1, $\Psi$ is clearly a non-zero preserving map. Moreover,
\begin{claim}\label{cl:psi-iff}
$J_{{\x}_k}(\T_k)\;=\; 0\; \iff\; \Psi(J_{{\x}_k}(\T_k))\;=\;0 $, over $\ringR[\x]$. 
\end{claim}

\begin{proof}
As $\pdeg(T_i) \le d$, each entry of the matrix can be of degree at most $d-1$; therefore $\pdeg(J_{{\x}_k}(\T_k)) \le  k(d-1)= D-1$. Thus, $\pdeg_{z}(\Psi(J_{{\x}_k}(\T_k))) <D$. Hence, the conclusion.
\end{proof}
\autoref{jacobian-main-eq} implies that 
\begin{align}
    \Psi(J_{{\x}_k}(\T_k))\;=\;\Psi(T_1 \cdots T_k) \cdot \sum_{g_1 \in L(T_1) , \hdots , g_k \in L(T_k)}\;\; \frac{ \Psi(J_{{\x}_k}(g_1,\hdots,g_k))}{\Psi(g_1 \hdots g_k)} \;. \label{jacobian-main-eq-psi}
\end{align}
As $T_i$ has product fanin $s$, the top-fanin in the sum in \autoref{jacobian-main-eq-psi} can be at most $s^k$. Then define, 
\begin{align}
\widetilde{F}\;:=\;\sum_{g_1 \in L(T_1) ,\hdots , g_k \in L(T_k)}\;\; \frac{\Psi(J_{{\x}_k}(g_1,\hdots,g_k))}{\Psi(g_1 \hdots g_k)} \;,\;\text{ over }\,\ringR[\x]. \label{eq:main-eq-F-psi}
\end{align}

\textbf{Well-definability of $\widetilde{F}$.} Note that, \[\Psi(g_i) \bmod z \ne 0 \implies 1/\Psi(g_1 \cdots g_k) \in \F[[\x,z]].\] Thus, RHS is an element in $\F[[\x,z]]$ and taking $\bmod\, z^D$ it is in $\ringR[\x]$. We remark that instead of minimally reducing $\bmod\,z^D$, we will work with an $F \in \F[z,\x]$ such that    $F=\tilde{F}$ over $\ringR[\x]$. Further, we ensure that the degree of $z$ is polynomially bounded.

\begin{claim}\label{cl:psi-to-F}
Over $\ringR[\x]$, $\Psi(J_{{\x}_k}(\T_k)) = 0  \iff F =0$.  
\end{claim}

\begin{proof}[Proof sketch]
This follows from the invertibility of $\Psi(T_1 \cdots T_k)$ in $R[\x]$.
%
\end{proof}

\textbf{The hitting set $H'$.} By $J_{\x_k}(\T_k) \ne 0$, and Claims \ref{cl:psi-iff}-\ref{cl:psi-to-F}, we have $F \ne 0$ over $\ringR[\x]$. We want to find $H' \subseteq \F^{n}$, s.t.~$\Psi(J_{\x_k}(\T_k))\rvert_{\x=\alphavec} \ne 0$, for some $\alphavec \in H'$ (which will ensure the rank-preservation). 
Towards this, we will show (below) that $F$ has $s^{O(\delta k)}$-size $\SWSP{}{}{[\delta]}$-circuit over $\ringR[\x]$. Next, \autoref{thm:swsp-hs} provides the hitting set $H'$ in time $s^{O(\delta^2 k \log s)}$. 


%

\begin{claim}[Main size bound]\label{claim:small-size-F-swsp}
$F \in \ringR[\x]$ has $\SWSP{}{}{[\delta]}$-circuit of size $ (s3^{\delta})^{O(k)}$.
\end{claim}

%
\noindent The proof studies the two parts of \autoref{eq:main-eq-F-psi}---
\begin{enumerate}
    \item The numerator $\Psi(J_{{\x}_k}(g_1,\hdots,g_k))$ has $O(3^{\delta}2^kk!ks)$-size $\SWSP{}{}{[\delta-1]}$-circuit (see \autoref{lem:num-size-thm2}), and
    \item $1/\Psi(g_1 \cdots g_k)$, for $g_i \in L(T_i)$ has $(s3^{\delta})^{O(k)}$-size $\SWSP{}{}{[\delta]}$-circuit; both over $\ringR[\x]$ (see \autoref{lem:denom-size-thm2}).
\end{enumerate}

\noindent We need the following two claims to prove the numerator size bound.

\begin{claim}\label{claim:jacobiansize}
    Let $g_i \in L(T_i)$, where $T_i \in \PSP{}{[\delta]}$ of size atmost $s$, then the polynomial $J_{{\x}_k}(g_1,\hdots,g_k)$ is computable by $\SPSP{[k!]}{[k]}{[\delta-1]}$ of size $O(k!\,ks)$.
\end{claim}
    
\begin{proof-sketch}
    Each entry of the matrix has degree at most $\delta-1$. Trivial expansion gives $k!$ top-fanin where each product (of fanin $k$) has size $\sum_{i}\,\size(g_i)$. As, $\size(T_i) \le s$, trivially each $\size(g_i) \le s$. Therefore, the total size is $k! \cdot\sum_{i}\,\size(g_i)=O(k!\,ks)$.
\end{proof-sketch}
    
\begin{claim} \label{claim:size-g}
    Let $g \in \Sigma \Pi^{\delta}$, then $\Psi(g) \in \Sigma \Pi^{\delta}$ of size $3^{\delta} \cdot \size(g)$ (for $n\gg\delta$).
\end{claim}
\begin{proof-sketch}
    Each monomial $\x^{\e}$ of degree $\delta$, can produce $\prod_i (e_i+1) \le ((\sum_{i} e_i +n)/n)^n \le (\delta/n+1)^n$-many monomials, by AM-GM inequality as $\sum_{i} e_i \le \delta$. As $\delta/ n \rightarrow 0$, we have $(1+\delta/n)^n \rightarrow e^{\delta}$. As $e < 3$, the upper bound follows.
\end{proof-sketch}

\begin{lemma}[Numerator size] \label{lem:num-size-thm2}
$\Psi(J_{{\x}_k}(g_1,\hdots,g_k))$ is computable by $\SWSP{}{}{[\delta-1]}$ of size $O(3^{\delta}\,2^kk\,k!s)=:s_2$.
\end{lemma}

\begin{proof}
    In \autoref{claim:jacobiansize} we showed that $J_{{\x}_k}(g_1,\hdots,g_k) \in \SPSP{[k!]}{[k]}{[\delta-1]}$ of size $O(k!ks)$. Moreover, for a $g \in \Sigma \Pi^{[\delta-1]}$, we have $\Psi(g) \in \Sigma \Pi^{[\delta-1]}$ of size at most $3^{\delta} \cdot \size(g)$, over $\ringR[\x]$ due to \autoref{claim:size-g}).
    
    Combining these, one concludes that $\Psi(J_{{\x}_k}(g_1,\hdots,g_k)) \in \SPSP{[k!]}{[k]}{[\delta-1]}$, of size $O(3^{\delta}\,k!ks)$. We \emph{convert} the $\Pi$-gate to $\wedge$ gate using waring identity (\autoref{lem:waring-rank}) which blowsup the size by a multiple of $2^{k-1}$. Thus, $\Psi(J_{{\x}_k}(g_1,\hdots,g_k)) \in \SWSP{}{}{[\delta-1]}$ of size $O(3^{\delta}\,2^kk\,k!s)$.
\end{proof}

In the following lemma, using power series expansion of expressions like $1/(1-a\cdot z)$, we conclude that $1/\Psi(g)$ has a small $\SWSP{}{}{[\delta]}$-circuit, which would further imply the same for $1/\Psi(g_1\cdots g_k)$.

\begin{lemma}[Denominator size] \label{lem:denom-size-thm2}
Let $g_i \in L(T_i)$. Then, $1/\Psi(g_1\cdots g_k)$ can be computed by a $\SWSP{}{}{[\delta]}$-circuit of size $s_1:=(s3^{\delta})^{O(k)}$, over $\ringR[\x]$.
\end{lemma}

\begin{proof}
Let $g \in L(T_i)$ for some $i$. Assume, $\Psi(g)=A-z\cdot B$, for some $A \in \F$ and $B \in \ringR[\x]$ of degree $\delta$, with $\size(B) \le 3^{\delta} \cdot s$, from \autoref{claim:size-g}. Note that, over $\ringR[\x]$,
\begin{align}
    \frac{1}{\Psi(g)} \;=\; \frac{1}{A(1- \frac{B}{A}\cdot z)}  \;=\; \frac{1}{A} \cdot \sum_{i=0}^{D-1} \left(\frac{B}{A}\right)^{i} \cdot z^i\;.\label{eq:powerseries-1/g-psi}
\end{align}
As, $B^i$ has a trivial $\WSP{}{[\delta]}$-circuit (over $\ringR[\x]$) of size $\le 3^{\delta}\cdot s + i$; summing over $i \in [D-1]$, the overall size is at most $D \cdot 3^{\delta} \cdot s + O(D^2)$. As $D < k \cdot d$, we conclude that $1/\Psi(g)$ has $\SWSP{}{}{[\delta]}$ of size $\poly(s \cdot k \cdot d 3^{\delta})$, over $\ringR[\x]$. Multiplying $k$-many such products directly gives an upper bound of $(s\cdot 3^{\delta})^{O(k)}$, using \autoref{lem:prod-swsw} (basically, waring identity).
\end{proof}

\begin{proof}[Proof of \autoref{claim:small-size-F-swsp}]
Combining Lemmas \ref{lem:num-size-thm2}-\ref{lem:denom-size-thm2}, observe that $\Psi(J_{\x_k}(\cdot)/\Psi(\cdot)$ has $\SWSP{}{}{[\delta]}$-circuit of size at most $(s_1 \cdot s_2)^2 = (s \cdot 3^{\delta})^{O(k)}$, over $\ringR[\x]$, using \autoref{lem:prod-swsw}. Summing up at most $s^k$ many terms (by defn.~of $F$), the size still remains $(s\cdot 3^{\delta})^{O(k)}$.
\end{proof}

{\bf \noindent Degree bound.} As, syntactic degree of $T_i$ are bounded by $d$, and $\Psi$ maintain $\pdeg_{\x}=\pdeg_{z}$, we must have $\pdeg_{z}(\Psi(J_{\x_k}(g_1,\hdots,g_k)) =\pdeg_{\x}(J_{\x_k}(g_1,\hdots,g_k)) \le D-1$. 
%
Note that, \autoref{lem:num-size-thm2} actually works over $\F[\x,z]$ and thus there is no additional degree-blow up (in $z$). However, there is some degree blowup in \autoref{lem:denom-size-thm2}, due to \autoref{eq:powerseries-1/g-psi}.

Note that \autoref{eq:powerseries-1/g-psi} shows that over $\ringR[\x]$,  \[\frac{1}{\Psi(g)} =\left( \frac{1}{A^D} \right)\,\cdot \left(\sum_{i=0}^{D-1} A^{D-1-i}z^{i} \cdot B^i\right) =: \frac{p(\x,z)}{q},\] where $q= A^{D}$. We think of $p \in \F[\x,z]$ and $q \in \F$. 
%
Note, $\pdeg_{z}(\Psi(g)) \le \delta$ implies $\pdeg_{z}(p) \le \pdeg_{z}((B\,z)^{D-1}) \le \delta \cdot (D-1)$. 
%

Finally, denote $1/\Psi(g_1 \cdots g_k)=: P_{g_1,\hdots,g_k}/Q_{g_1,\hdots,g_k}$, over $\ringR[\x]$. This is just multiplying $k$-many $(p/q)$'s; implying a degree blowup by a multiple of $k$. In particular -- $\pdeg_{z}(P_{(\cdot)}) \le \delta \cdot k \cdot (D-1)$
%
%
%
%
%
%
Thus, in \autoref{eq:main-eq-F-psi}, summing up $s^k$-many terms gives an expression (over $\ringR[\x]$):
\[ F \;=\; \sum_{g_1 \in L(T_1),\hdots, g_k \in L(T_k)}\; \Psi(J_{\x_k}(g_1,\hdots,g_k)) \cdot \left( \frac{P_{g_1,\hdots,g_k}}{Q_{g_1,\hdots,g_k}} \right)\;=:\; \frac{P(\x,z)}{Q}\;.\]

Verify that $Q \in \F$.
%
%
The degree of $z$ also remains bounded by \[\max_{g_i \in L(T_i), i\in[k]} \pdeg_{z}(P_{g_1,\hdots,g_k})+\delta k \le \poly(s).\]
%
Using the degree bounds, we finally have $P \in \F[\x,z]$ as a $\SWSP{}{}{[\delta]}$-circuit (over $\F(z)$) of size $n^{O(\delta)}\,(s3^{\delta})^{O(k)}$ $=3^{O(\delta k)} s^{O(k+\delta)} =: s_3$. 

We want to \emph{construct} a set $H' \subseteq \F^n$ such that the action $P(H',z) \ne 0$.  Using \cite{forbes2015deterministic} (\autoref{thm:swsp-hs}), we conclude that it has $s^{O(\delta\log s_3)}$ $= s^{O(\delta^2 k \log s)}$ size hitting set which is constructible in a similar time. 
Hence, the construction of $\Phi$ follows, making $\Phi(f)$ a $k+2$ variate polynomial. Finally, by the obvious degree bounds of $\y,z,t$ from the definition of $\Phi$, we get the blackbox PIT algorithm with time-complexity $s^{O(\delta^2 k \log s)}$; finishing \autoref{thm:thm2}(b).

We could also give the final hitting set for the general problem.

\begin{proof}[Solution to \autoref{prob:prob1-thm2}]
We know that \[C(T_1,\hdots,T_m) = 0 \iff E:= \Phi(C(T_1,\hdots,T_m)) =0.\] 
 Since, $H'$ can be constructed in $s^{O(\delta^2\,k\,\log s)}$-time, it is trivial to find hitting set for $E\rvert_{H'}$ (which is just a $k+2$-variate polynomial with the aformentioned degree bounds). The final hitting set for $E$ can be constructed in $s'^{O(k)} \cdot s^{O(\delta^2\,k\,\log s)}$-time.
\end{proof}

%

\begin{remark*}
\begin{enumerate}
    \setlength\itemsep{1em}
    \item As Jacobian Criterion (\autoref{fact:jacobian-criterion}) holds when the characteristic is $> d^{\,\tdeg}$, it is easy to conclude that our theorem holds for all fields of char $> d^k$.
    \item The above proof gives an efficient reduction from blackbox PIT for $\SPSP{[k]}{}{[\delta]}$ circuits to $\SWSP{}{}{[\delta]}$ circuits. In particular, a poly-time hitting set for $\SWSP{}{}{[\delta]}$ circuits would put PIT for $\SPSP{[k]}{}{[\delta]}$ in $\P$. 
    \item Also, $\DiDI$-technique (of Theorem \autoref{thm:thm1}) directly gives a blackbox algorithm, but the complexity is {\em exponentially} worse (in terms of $k$ in the exponent) for its recursive blowups.
\end{enumerate}
\end{remark*}

\subsection{PIT for \texorpdfstring{$\SPSU{[k]}{}{}$}{SkPSU}}

As we remarked earlier, the proof of \autoref{thm:thm2}(a) is similar to the one we discussed in section \ref{sec:SkPSPdelta}. Here we sketch the proof, stating some relevant changes. Similar to \autoref{thm:thm2}(b), we generalize this theorem and prove for a much bigger class of polynomials.  

\begin{problem}\label{prob:prob2-thm2}
 Let $\{T_i \,|\, i \in [m]\}$ be $\PSU{}{}$ circuits of (syntactic) degree at most $d$ and size $s$. Let the transcendence degree of $T_i$'s, $\tdeg_{\F}(T_1,\hdots,T_m) =:k \ll s$. Further, $C(x_1,\hdots,x_m)$ be a circuit of size + degree $<s'$. Design a blackbox-PIT algorithm for $C(T_1,\hdots,T_m)$.
\end{problem}

It is trivial to see that $\SPSU{[k]}{}{}$ is a very special case of the above settings. We will use the same idea (\& notation) as in \autoref{thm:thm2}(b), using the Jacobian technique. The main idea is to come up with $\Psi$ map, and correspondingly the hitting set $H'$. If $g \in L(T_i)$, then $\size(g) \le O(dn)$. 
%
The $D$ (and hence $R[\x]$) remains as before. Claims \ref{cl:psi-iff}-\ref{cl:psi-to-F} hold similarly. We will construct the hitting set $H'$ by showing that $F$ has a small $\SWSU{}{}{}$ circuit over $R[\x]$. 

Note that, \autoref{claim:jacobiansize} remains the same for $\SWSU{}{}{}$ (implying the same size blowup). However, \autoref{claim:size-g}, the size blowup is $O(d\,\size(g))$, because each monomial $x^e$ can only produce $d+1$ many monomials. Therefore, similar to \autoref{lem:denom-size-thm2}, one can show that $\Psi(J_{\x_k}(g_1,\hdots,g_k)) \in \SWSU{}{}{}$, of size $O(2^kk!kds)$. Similarly, the size in \autoref{lem:num-size-thm2} can be replaced by $s^{O(k)}$. Therefore, we get (similar to \autoref{claim:small-size-F-swsp}):

\begin{claim}
$F \in R[\x]$ has $\SWSU{}{}{}$-circuit of size $s^{O(k)}$.
\end{claim}

Next, the degree bound also remains the same. Following the same footsteps, it is not hard to see that 
%
while degree bound on $z$ remains $\poly(ksd)$. Therefore, $P \in \F[\x,z]$ has $\SWSU{}{}{}$-circuit of size $s^{O(k)}$.

We want to \emph{construct} a set $H' \subseteq \F^n$ such that the action $P(H',z) \ne 0$.  By \autoref{lem:swsu-hs}, we conclude that it has $s^{O(k\,\log \log s)}$ size hitting set  which is constructible in a similar time. Hence, the construction of map $\Phi$ and the theorem follows (from $z$-degree bound).


\begin{proof}[Solution to \autoref{prob:prob2-thm2}]
We know that \[C(T_1,\hdots,T_m) = 0 \iff E:= \Phi(C(T_1,\hdots,T_m)) =0.\] 
Since, $H'$ can be constructed in $s^{O(k \log \log s)}$ time, it is trivial to find hitting set for $E\rvert_{H'}$ (which is just a $k+2$-variate polynomial with the aforementioned degree bounds). The final hitting set for $E$ can be constructed in $s'^{O(k)} \cdot s^{O(k \log \log s)}$ time.
\end{proof}
